% sections/lagrangian.tex

\section{Identification with exceptional supergravity}
\label{sec:identification}

The algebraic data of $\hthree$ (the Peirce decomposition, the cubic
norm $\det(X)$, and the $\Esix$ structure group) uniquely match the
exceptional entry in the Gunaydin-Sierra-Townsend (GST) classification
of $N{=}2$ Maxwell-Einstein supergravity theories. This identification
plays the same role as the Jordan-von~Neumann-Wigner classification in
Paper~5: the algebraic structure is proved, and a classification theorem
identifies it with a known physical theory. The $N{=}2$ supersymmetric
structure is \emph{identified} via the GST bijection as the unique match
to the algebraic data, not assumed as an input.

\subsection{The GST classification}

Gunaydin, Sierra, and Townsend~\cite{GST1983,GST1984} (see
also~\cite{Gunaydin2008} for a modern review) proved that
degree-3 Euclidean Jordan algebras with positive-definite trace form are
in bijection with $N{=}2$ Maxwell-Einstein supergravity theories in five
dimensions. Given such an algebra $J$ with cubic norm $N(h)$, the
5-dimensional bosonic Lagrangian is uniquely determined:
\begin{align}
\label{eq:5d-lagrangian}
e^{-1}\mathcal{L}_5 &= -\frac{R}{2}
  - g_{ij}\,\partial_\mu \phi^i\,\partial^\mu \phi^j
  - \tfrac{1}{4}\,G_{IJ}\,F^I_{\mu\nu}\,F^{J\mu\nu}
  \nonumber\\
  &\quad - \tfrac{C_{IJK}}{24}\,\epsilon^{\mu\nu\rho\sigma\tau}\,
    F^I_{\mu\nu}\,F^J_{\rho\sigma}\,A^K_\tau,
\end{align}
where $C_{IJK} = \frac{1}{6}\,d_{IJK}$, the scalar metric $g_{ij}$
and gauge kinetic matrix $G_{IJ}$ are derived from the cubic norm $N$
via very special real (VSR) geometry, and $R$ is the Ricci scalar. The
Einstein-Hilbert term $-R/2$ is part of the GST output: the bijection
identifies the full bosonic Lagrangian, including gravity.

\begin{proposition}
\label{prop:gst-hypotheses}
The algebra $\hthree$ satisfies all three GST hypotheses:
\begin{enumerate}[label=(H\arabic*)]
\item \textbf{Degree 3:} $\hthree$ has a cubic norm $\det(X)$
  (Theorem~\ref{thm:springer}).
\item \textbf{Formally real:} $\hthree$ is one of the five algebras
  in the Jordan-von~Neumann-Wigner
  classification~\cite{JordanVonNeumannWigner1934} that are formally
  real (equivalently: $\jp{X}{X} = 0$ implies $X = 0$).
\item \textbf{Positive-definite trace form:} The bilinear form
  $\Tr(\jp{X}{X}) = \alpha^2 + \beta^2 + \gamma^2 + 2(|x_1|^2 +
  |x_2|^2 + |x_3|^2) > 0$ for $X \neq 0$. On the 26-dimensional
  tangent space at the identity, all eigenvalues are positive
  (minimum eigenvalue $\frac{1}{4}$).
\end{enumerate}
\end{proposition}

By the GST bijection, $\hthree$ therefore uniquely determines a
5-dimensional $N{=}2$ MESGT: the exceptional magic supergravity.
The 4-dimensional theory is obtained via the $r$-map (dimensional
reduction on a circle)~\cite{deWitVanProeyen1992}, giving a special
K\"ahler scalar manifold
\begin{equation}
\label{eq:scalar-manifold}
\mathcal{M} = \frac{\Eseven}{E_{6(-78)} \times \UU{1}},
\quad \dim_{\mathbb{R}} = 54.
\end{equation}

\subsection{Field content}

The 4-dimensional field content, organized by Peirce sector, is:

\begin{center}
\begin{tabular}{lccc}
\toprule
Peirce sector & $I$ range & Vectors & Scalars \\
\midrule
$\Vone$ (observer) & $1$ & 1 & 2 \\
$\Vhalf$ (SM matter) & $2$--$17$ & 16 & 32 \\
$\Vzero$ (spacetime + internal) & $18$--$27$ & 10 & 20 \\
Graviphoton & $0$ & 1 & 0 \\
\midrule
Total & & 28 & 54 \\
\bottomrule
\end{tabular}
\end{center}

The 28 vectors include the graviphoton ($I = 0$) from the gravity
multiplet and 27 matter vectors from the Peirce decomposition.
The 54 real scalars parametrize the K\"ahler manifold
$\mathcal{M}$~\cite{FerraraGunaydin2006}. After dimensional reduction
via the $r$-map, all 27 Jordan algebra directions contribute to these
scalars, including the 6 internal $W$-directions from
$\ker(\pi_u)$ and the 4 spacetime-identified directions of $\Vzero$.
The $W$-directions are not lost in the reduction; they become internal
scalar moduli of the 4-dimensional theory.

\subsection{Coupling decomposition}

The $C_{IJK} = \frac{1}{6}\,d_{IJK}$ tensor governs all
matter-gravity couplings. Its 106 nonzero entries decompose into three
physical channels:

\begin{enumerate}
\item \textbf{Gravitational self-coupling} ($\Vone, \Vzero, \Vzero$):
  10 entries, reproducing the $\det_2$ bilinear form on $\Vzero$. This
  couples the observer degree of freedom ($I = 1$) to the spacetime
  sector.

\item \textbf{Matter-spacetime coupling} ($\Vhalf, \Vhalf, \Vzero$
  restricted to the 4 spacetime directions of $\htwo{C}_u$): 48 entries.
  Per-index counts: indices 18 and 19 each couple to 16 matter fields;
  indices 20 and 27 each couple to 8.

\item \textbf{Matter-internal coupling} ($\Vhalf, \Vhalf, \Vzero$
  restricted to the 6 internal $W$-directions): 48 entries, 8 per
  internal index.
\end{enumerate}

The coupling is \emph{universal}: every matter field in $\Vhalf$ couples
to every spacetime direction. This universality is a structural property
of the $d_{IJK}$ tensor, reflecting the algebraic coupling of all matter
to the spacetime sector.

\subsection{The 4-dimensional Lagrangian}

The complete 4-dimensional bosonic Lagrangian, identified via the GST
bijection and $r$-map, is~\cite{deWitVanProeyen1992,
LauriaVanProeyen2020}:
\begin{align}
\label{eq:4d-lagrangian}
e^{-1}\mathcal{L}_{\mathrm{bos}} &= -\frac{R}{2}
  + g_{i\bar{j}}\,\partial_\mu z^i \partial^\mu \bar{z}^{\bar{j}}
  \nonumber\\
  &\quad + \im(\mathcal{N}_{IJ})\, F^I_{\mu\nu} F^{J\mu\nu}
  + \re(\mathcal{N}_{IJ})\, F^I_{\mu\nu} {*F}^{J\mu\nu},
\end{align}
where:
\begin{itemize}
\item $g_{i\bar{j}} = \partial_i \partial_{\bar{j}} K$ is the
  K\"ahler metric, with K\"ahler potential $K$ derived from the
  prepotential $F(X) = d_{IJK} X^I X^J X^K / (6 X^0)$, where
  $I, J, K = 1, \ldots, 27$ run over the matter fields and $X^0$ is
  the projective coordinate (graviphoton section).

\item $\mathcal{N}_{IJ}$ is the gauge kinetic matrix (period matrix),
  also derived from $F(X)$ via special K\"ahler geometry.

\item $R$ is the 4-dimensional Ricci scalar.
\end{itemize}

\begin{remark}[What $\det(X)$ determines]
\label{rem:precise-claim}
The prepotential $\det(X)$ determines the scalar manifold
$\mathcal{M}$, the gauge kinetic matrix $\mathcal{N}_{IJ}$, and all
matter-gravity couplings $C_{IJK}$. The Einstein-Hilbert term $-R/2$
is part of the identified MESGT Lagrangian: the GST bijection maps the
algebraic data of $\hthree$ to the \emph{complete} bosonic Lagrangian
including the gravitational sector. The identification is with a theory
that has gravity; gravity is not independently derived from $\det(X)$.
\end{remark}

\subsection{$N{=}2$ supersymmetry as a consequence}
\label{sec:n2-derived}

In the construction above, $N{=}2$ supersymmetry was not assumed.
The bosonic Lagrangian is established in two logically independent
steps: very special real (VSR) geometry gives the Lagrangian from the
cubic prepotential without any supersymmetric input, and the GST
classification then identifies the same Lagrangian as the bosonic
sector of the exceptional $N{=}2$ MESGT.

\paragraph{Step 1: VSR geometry (no SUSY input).}
The chain from $\hthree$ to the bosonic Lagrangian is:

\begin{enumerate}[label=(\roman*)]
\item $\hthree$ is a degree-3 Euclidean Jordan algebra (algebraic
  definition).
\item $\Ffour = \Aut(\hthree)$ gives the $\rep{27}$ representation
  (representation theory).
\item $\Esix$ structure group defines the scalar manifold (symmetric
  space theory).
\item The VSR metric $g_{ij}$ is $\Esix$-covariant (Schur's lemma).
\item The cubic prepotential $V = C_{IJK} h^I h^J h^K$ is unique
  (Springer).
\item The gauge kinetic matrix $G_{IJ}$ is uniquely fixed (VSR
  geometry~\cite{deWitVanProeyen1992}).
\item Gauge invariance determines the Chern-Simons coupling
  (gauge theory).
\end{enumerate}

Steps (i)--(vii) produce the unique bosonic Lagrangian determined by
the cubic prepotential $\det(X)$ within VSR geometry. No
supersymmetry appears at any step. (VSR geometry was developed within
the $N{=}2$ supergravity program~\cite{deWitVanProeyen1992}; the
steps above use only the cubic-polynomial structure, not SUSY
variations, but the framework's origin should be noted.)

\paragraph{Step 2: GST classification ($N{=}2$ identification).}
The GST classification~\cite{GST1983,GST1984} is exhaustive within its
category: every degree-3 formally real Jordan algebra with
positive-definite trace form corresponds to a unique $N{=}2$ MESGT. The
existence of this bijection was established by GST through explicit
verification of $N{=}2$ supersymmetry closure. Since $\hthree$ satisfies
all three GST hypotheses (Proposition~\ref{prop:gst-hypotheses}), the
GST bijection identifies the VSR Lagrangian from Step~1 as the bosonic
sector of the exceptional $N{=}2$ MESGT.

The $N{=}2$ label appears only in Step~2, as the \emph{output} of
the classification, not as an input to any step that builds the
Lagrangian. A circularity audit confirms that supersymmetry is absent
from Steps (i)--(vii). To be precise: the $N{=}2$ label enters because
the GST classification theorem classifies $N{=}2$ MESGTs; it is not
derived from the self-modeling framework independently of this
classification context. The VSR and GST routes give the same Lagrangian,
confirming that the result is not an artifact of the $N{=}2$ framework.

\subsection{Cosmological constant}

The ungauged $N{=}2$ MESGT has an identically vanishing scalar potential
$V(z, \bar{z}) = 0$~\cite{LauriaVanProeyen2020}, giving $\Lambda = 0$
at the classical level. The observed cosmological constant
$\Lambda > 0$ is not explained by this framework. Obtaining $\Lambda
\neq 0$ would require either gauging the theory (introducing a scalar
potential) or invoking quantum corrections. This is an open limitation,
shared with essentially all approaches to the cosmological constant
problem.
